\subsection{Latar Belakang}

Perkembangan teknologi informasi dan komunikasi dalam beberapa tahun terakhir telah mendorong transformasi signifikan dalam penyediaan layanan psikologi, khususnya melalui pemanfaatan aplikasi mobile sebagai media penyajian tes psikologi. Aplikasi berbasis mobile memungkinkan pengguna untuk mengakses berbagai jenis tes psikologi secara fleksibel, cepat, dan efisien, sehingga mendukung proses pemahaman diri dan pengembangan potensi individu. Tes psikologi seperti tes kecerdasan, tes kecemasan, tes stres, tes bakat, tes potensi, dan tes ketangguhan merupakan instrumen penting dalam mengukur aspek kognitif dan non-kognitif individu. Dalam kajian psikometri modern, kualitas hasil pengukuran sangat ditentukan oleh validitas dan reliabilitas instrumen yang digunakan. Validitas menekankan pada ketepatan alat ukur dalam merepresentasikan konstruk psikologis yang dinilai, sedangkan reliabilitas menunjukkan konsistensi hasil pengukuran dalam kondisi yang serupa (Flake et al., 2019). Pendekatan kontemporer dalam psikometri juga menekankan pentingnya pemahaman hubungan antara konstruk teoretis dan skor tes untuk memastikan interpretasi hasil yang akurat dan bermakna (Zumbo \& Chan, 2021).

Aplikasi Jatidiri dikembangkan sebagai aplikasi mobile yang menyediakan berbagai tes psikologi untuk membantu pengguna mengenali kondisi psikologis dan potensi diri mereka. Penyajian tes psikologi dalam bentuk digital sejalan dengan tren asesmen psikologis berbasis teknologi yang menekankan aksesibilitas, efisiensi waktu, dan kemudahan penggunaan. Penelitian terbaru menunjukkan bahwa asesmen psikologi berbasis digital dapat memberikan hasil yang sebanding dengan metode konvensional apabila dirancang dengan struktur pengujian yang jelas serta antarmuka pengguna yang baik (Coyne et al., 2020). Oleh karena itu, selain aspek konseptual psikometri, kualitas desain antarmuka pengguna dan pengalaman pengguna menjadi faktor krusial dalam memastikan efektivitas aplikasi psikologi berbasis mobile.

Proses pengembangan aplikasi Jatidiri difokuskan pada implementasi desain antarmuka pengguna yang telah dirancang menggunakan Figma ke dalam aplikasi mobile berbasis Flutter. Dalam pengembangan perangkat lunak modern, metode Prototype merupakan pendekatan yang efektif untuk membangun aplikasi secara iteratif melalui pembuatan rancangan awal, evaluasi, dan penyempurnaan berkelanjutan berdasarkan umpan balik pengguna atau pemangku kepentingan. Metode ini sangat sesuai untuk pengembangan antarmuka aplikasi karena memungkinkan penyesuaian desain secara bertahap hingga tercapai kesesuaian antara rancangan visual dan implementasi teknis (Santos et al., 2020). Pendekatan prototyping juga mendukung pengembangan aplikasi yang lebih berorientasi pada pengguna, terutama dalam konteks aplikasi psikologi yang membutuhkan kejelasan navigasi dan kemudahan interaksi (Preece et al., 2021).

Pemilihan framework Flutter dalam pengembangan aplikasi Jatidiri didasarkan pada kemampuannya dalam menghasilkan aplikasi mobile lintas platform dengan performa yang optimal dan tampilan antarmuka yang konsisten. Flutter menyediakan berbagai komponen antarmuka yang memungkinkan penerjemahan desain UI/UX dari Figma ke dalam kode secara efisien dan terstruktur. Penelitian terkini menunjukkan bahwa penggunaan framework lintas platform seperti Flutter dapat meningkatkan produktivitas pengembangan serta menjaga konsistensi antarmuka pengguna pada berbagai perangkat (Alshamrani \& Bahattab, 2022). Dengan mengintegrasikan metode Prototype dan framework Flutter, pengembangan aplikasi Jatidiri diharapkan mampu menghasilkan antarmuka aplikasi yang sesuai dengan desain awal, responsif, dan mudah digunakan oleh pengguna dalam mengakses berbagai tes psikologi.

\subsection{Identifikasi Masalah}

1. Bagaimana proses perancangan dan implementasi aplikasi Jatidiri berbasis mobile menggunakan framework Flutter dengan menerapkan metode Prototype agar sesuai dengan kebutuhan pengguna dan rancangan sistem?

2. Bagaimana cara menerjemahkan desain antarmuka pengguna (UI) dari Figma ke dalam komponen Flutter sehingga menghasilkan tampilan yang interaktif, responsif, dan mudah digunakan?

3. Bagaimana memastikan bahwa fungsi-fungsi utama aplikasi, khususnya penyajian berbagai tes psikologi dan halaman komunitas, dapat berjalan sesuai dengan kebutuhan dan rancangan yang telah ditetapkan

\subsection{Tujuan dan Manfaat}

\subsubsection{Tujuan}

Adapun tujuan dari penelitian ini, tujuannya sebagai berikut:

1. Merancang dan mengimplementasikan aplikasi Jatidiri berbasis mobile menggunakan framework Flutter dengan menerapkan metode Prototype.

2. Mengimplementasikan desain antarmuka pengguna (UI) dari Figma ke dalam aplikasi Flutter agar menghasilkan tampilan yang interaktif, responsif, dan mudah digunakan.

3. Menghasilkan aplikasi mobile yang mampu menyajikan berbagai tes psikologi serta memastikan fungsi dan antarmuka aplikasi berjalan sesuai dengan kebutuhan dan rancangan yang telah ditetapkan.

\subsubsection{Manfaat}

Adapun Manfaat dari penelitian ini, manfaatnya sebagai berikut:

1. Menambah referensi ilmiah dalam bidang pengembangan aplikasi mobile, khususnya terkait implementasi desain UI/UX ke dalam aplikasi berbasis Flutter.

2. Memberikan gambaran penerapan metode Prototype dalam pengembangan antarmuka aplikasi mobile yang berorientasi pada pengguna.

3. Memberikan kontribusi dalam pengembangan aplikasi tes psikologi berbasis teknologi mobile dengan memanfaatkan framework modern seperti Flutter.

\subsection{Ruang Lingkup}

Ruang lingkup penelitian mencakup sebagai berikut:

a. Implementasi desain antarmuka pengguna (UI/UX) aplikasi Jatidiri yang responsif dan ramah pengguna berdasarkan desain yang dibuat menggunakan Figma.

b. Pengembangan antarmuka aplikasi mobile menggunakan framework Flutter dan bahasa pemrograman Dart untuk menyajikan berbagai fitur tes psikologi serta halaman komunitas.

c. Penerapan metode Prototype dalam proses pengembangan aplikasi, yang meliputi pembuatan rancangan awal, implementasi, evaluasi, dan penyempurnaan antarmuka aplikasi.

d. Pengujian aplikasi untuk memastikan setiap fitur dan tampilan antarmuka berjalan sesuai dengan kebutuhan dan desain yang telah ditentukan menggunakan metode black box testing.

e. Perangkat lunak yang digunakan dalam pengembangan aplikasi meliputi Flutter SDK, Android Studio atau Visual Studio Code, Figma, serta emulator atau perangkat Android untuk proses pengujian.

\subsection{Rencana Kegiatan}

\begin{table}[h!]
\centering
\caption{Rencana Kegiatan Periode Oktober - Desember}
\label{tab:rencana_kegiatan}
\small
\begin{tabular}{|c|l|c|c|c|c|c|c|c|c|c|c|c|}
\hline
\rowcolor{yellow!60}
\multirow{2}{*}{No.} & \multirow{2}{*}{Kegiatan} & \multicolumn{12}{c|}{Bulan} \\
\cline{3-14}
\rowcolor{yellow!60}
 &  & \multicolumn{4}{c|}{Oktober} & \multicolumn{4}{c|}{November} & \multicolumn{4}{c|}{Desember} \\
\cline{3-14}
\rowcolor{yellow!60}
 &  & 1 & 2 & 3 & 4 & 1 & 2 & 3 & 4 & 1 & 2 & 3 & 4 \\
\hline

1 & Orientasi dan Pengenalan Proyek 
& \cellcolor{gray!60} &  &  &  &  &  &  &  &  &  &  &  \\
\hline

2 & Analisis Desain Figma dan User Flow 
&  & \cellcolor{gray!60} & \cellcolor{gray!60} &  &  &  &  &  &  &  &  &  \\
\hline

3 & Implementasi Halaman Awal 
&  &  &  & \cellcolor{gray!60} &  &  &  &  &  &  &  &  \\
\hline

4 & Implementasi Halaman Autentikasi 
&  &  &  &  & \cellcolor{gray!60} &  &  &  &  &  &  &  \\
\hline

5 & Implementasi Halaman Survei dan Home 
&  &  &  &  &  & \cellcolor{gray!60} &  &  &  &  &  &  \\
\hline

6 & Implementasi Halaman Tes 
&  &  &  &  &  &  & \cellcolor{gray!60} & \cellcolor{gray!60} &  &  &  &  \\
\hline

7 & Implementasi Setiap Halaman Macam - Macam Tes 
&  &  &  &  &  &  &  &  & \cellcolor{gray!60} & \cellcolor{gray!60} &  &  \\
\hline

8 & Implementasi Halaman Komunitas 
&  &  &  &  &  &  &  &  &  &  & \cellcolor{gray!60} &  \\
\hline

9 & Pengujian Black Box dan penyusunan laporan 
&  &  &  &  &  &  &  &  &  &  &  & \cellcolor{gray!60} \\
\hline

\end{tabular}
\end{table}

